\documentclass{article}

% NeurIPS 2026 Evaluations & Datasets Track --- single-blind opt-in
% (allowed for dataset-centered submissions where the dataset organization
%  and prior peer-reviewed work cannot be released anonymously without
%  breaking provenance, e.g. JuDDGES on Hugging Face).
\usepackage[eandd, nonanonymous]{neurips_2026}
% TODO: After acceptance, switch to camera-ready by adding `final`:
%   \usepackage[eandd, nonanonymous, final]{neurips_2026}
% Preprint variant for arXiv:
%   \usepackage[preprint]{neurips_2026}
%
% TODO(reviewer-facing): if NeurIPS 2026 publishes an updated style file,
% replace neurips_2026.sty in this directory before submission.

\usepackage[utf8]{inputenc}
\usepackage[T1]{fontenc}
\usepackage[main=english,polish]{babel} % Polish hyphenation for embedded Polish content
\usepackage{hyperref}
\usepackage{url}
\usepackage{booktabs}
\usepackage{tabularx}
\usepackage{graphicx}
\usepackage{amsmath}
\usepackage{amsfonts}
\usepackage{nicefrac}
\usepackage{microtype}
\usepackage{xcolor}
\usepackage{fancyvrb}
\usepackage[most]{tcolorbox}
\usepackage{natbib}
\bibliographystyle{plainnat}

% Reusable styled boxes for sample judgment excerpts and extraction-prompt
% listings in the appendix. Keep both variants: examplebox (non-breakable)
% and promptbox (breakable across pages).
\newtcolorbox{examplebox}[1]{%
  enhanced,
  colback=gray!4, colframe=black!55, boxrule=0.4pt, arc=2pt,
  title={#1}, fonttitle=\sffamily\bfseries\small,
  coltitle=white, colbacktitle=black!65,
  left=4pt, right=4pt, top=3pt, bottom=3pt,
  before skip=6pt, after skip=6pt,
}
\newtcolorbox{promptbox}[1]{%
  enhanced, breakable,
  colback=gray!4, colframe=black!55, boxrule=0.4pt, arc=2pt,
  title={#1}, fonttitle=\sffamily\bfseries\small,
  coltitle=white, colbacktitle=black!65,
  left=4pt, right=4pt, top=3pt, bottom=3pt,
  before skip=6pt, after skip=6pt,
}
% Silver-label callout used for sample silver-label excerpts in section
% bodies. Visual tier between examplebox (judgment quotations) and
% promptbox (extraction prompts) -- lighter frame to signal "machine
% generated, not gold".
\newtcolorbox{silverbox}[1]{%
  enhanced, breakable,
  colback=gray!2, colframe=black!40, boxrule=0.4pt, arc=2pt,
  title={#1}, fonttitle=\sffamily\bfseries\small,
  coltitle=white, colbacktitle=black!50,
  left=4pt, right=4pt, top=3pt, bottom=3pt,
  before skip=6pt, after skip=6pt,
}

% Convenience macros (mirrors shared/settings.tex from JuDDGES 2024)
\newcommand{\juddges}{JuDDGES}
\newcommand{\juddgespl}{JuDDGES-pl}

\title{JuDDGES-pl: Hundreds of Thousands of Structurally Enriched Polish Judgments from Common and Administrative Courts}

% Authors match the OpenReview submission (single-blind, allowed for
% dataset-centered E&D submissions). Author list is final; the broader
% JuDDGES contributor list in CITATION.cff is intentionally not used
% for this submission, which is scoped to the four WUST authors who
% led the corpus release described in this paper.
\author{%
  {\L}ukasz Augustyniak \\
  Wroclaw University of \\ Science and Technology \\ (WUST) \\
  \And
  Jakub Binkowski \\
  WUST \\
  \And
  Albert Sawczyn \\
  WUST \\
  \And
  Tomasz Kajdanowicz \\
  WUST \\
}

\begin{document}

\maketitle

\begin{abstract}
Civil-law jurisdictions remain severely underrepresented in legal NLP resources, where existing corpora are dominated by English-language common-law data and rarely include analytical structure beyond raw text. We release \juddgespl{} (\url{https://huggingface.co/datasets/JuDDGES/juddges-pl}), a single large-scale, analytically enriched Polish court judgment corpus spanning the two disjoint branches of the Polish judiciary: a common-courts subset covering the Supreme Court, courts of appeal, regional and district courts, and an administrative-courts subset covering the Supreme Administrative Court and lower administrative courts. The corpus comprises 2{,}691{,}842 judgments spanning 1981--2025 (437{,}450 common-court judgments from 1994--2025 and 2{,}254{,}392 administrative-court judgments from 1981--2025), each augmented with sixteen LLM-extracted analytical fields --- including structured factual state, legal state, summaries, theses, statutory bases, legal concepts, parties, outcomes, and keywords --- generated through a uniform Google Gemini 2.5 Pro extraction pipeline, alongside court-hierarchy, chamber, and procedural metadata preserved from the source portals.

Our contribution is curatorial and infrastructural. We describe the acquisition, normalization, and large-scale LLM-based enrichment pipeline that converts heterogeneous court-publication formats into a uniform, analysis-ready resource; we characterize the corpus through descriptive analyses of coverage, document structure, citation patterns, court hierarchy, extraction-field statistics, and pseudonymization fidelity; and we document the resource using Croissant metadata with Responsible AI fields and an accompanying evaluation card. \textbf{We treat LLM-extracted fields explicitly as silver labels} --- useful for analysis, model pretraining, and weak-supervision pipelines, but \textbf{not as gold annotations for benchmarking} --- and document this distinction throughout.

\juddgespl{} is designed to enable downstream evaluation of civil-law legal NLP --- including structured information extraction, statute-grounded reasoning, cross-branch generalization between common and administrative courts, longitudinal analysis of judicial language, and fairness audits in jurisdictions outside the Anglophone common-law mainstream --- without itself proposing new tasks or model comparisons. The corpus is openly released on Hugging Face under CC BY 4.0, alongside the reproducible enrichment pipeline (Apache 2.0) and validated Croissant metadata.
\end{abstract}

% Introduction
% 01-introduction.tex
% Goal: establish (i) the civil-law gap, (ii) what JuDDGES-pl is,
% (iii) what counts as our contribution, (iv) the silver-label framing
% as an explicit design commitment.

\section{Introduction}
\label{sec:introduction}

\subsection{The civil-law gap in legal NLP}
\label{sec:intro:gap}

Contemporary evaluation in legal NLP is dominated by suites such as LegalBench~\citep{guhaLegalbenchCollaborativelyBuilt2023a}, LexGLUE~\citep{chalkidis-etal-2022-lexglue}, and CaseHOLD~\citep{zheng2021does}, all of which are English-language and constructed almost exclusively from common-law sources. Civil-law systems, however, differ from their common-law counterparts in ways that are not merely stylistic: judgments cite codified statutes rather than precedential cases, document structure follows a fixed two-part form (\emph{sentencja} and \emph{uzasadnienie}), and the courts themselves are organized into separate hierarchies---common and administrative---governed by disjoint procedural codes. These structural features shape what models must learn in order to reason about civil-law texts, and they are largely absent from the resources on which legal language models are currently trained and evaluated. Polish legal NLP has produced valuable building blocks, including PolEval campaigns, the Bielik training corpora, and the original JuDDGES release at NeurIPS 2024~\citep{augustyniak2024juddges}, but no existing resource simultaneously covers two judicial branches in a unified release, supplies an LLM-derived analytical layer at scale, and is documented to the standards expected of a NeurIPS Datasets \& Benchmarks contribution. The result is a persistent gap between the empirical centre of legal NLP and the legal systems that govern most of continental Europe. \juddgespl{} is designed to narrow that gap for Polish, while making the methodology transparent enough to be reused for other civil-law jurisdictions.

\subsection{What we release}
\label{sec:intro:release}

\juddgespl{} is a unified Hugging Face dataset, distributed under the identifier \texttt{JuDDGES/juddges-pl}, that consolidates two complementary configurations: \texttt{pl-court}, covering Polish common courts, and \texttt{pl-nsa}, covering the administrative court branch. The release contains hundreds of thousands of judgments spanning multiple years of jurisprudence, and the underlying texts are preserved verbatim, with the anonymization applied by the publishing courts left intact rather than re-processed. This design choice keeps the corpus faithful to the documents that practitioners and researchers actually encounter, and it ensures that any subsequent analysis of pseudonymization fidelity is performed on the same surface form that downstream users will see.

Layered on top of these documents is an LLM-extracted analytical layer produced uniformly with Google Gemini 2.5 Pro under a single, version-pinned prompting protocol. For every judgment we extract a factual state, a legal state, summaries, theses, statutory bases, and keywords, providing a structured view of the case alongside the raw text. The release is accompanied by validated Croissant metadata with full Responsible AI fields, and the entire extraction pipeline---prompts, schemas, model version, and post-processing---is shipped in reproducible form so that the analytical layer can be regenerated, audited, or extended without recourse to undocumented tooling.

\subsection{Contributions}
\label{sec:intro:contributions}

\begin{enumerate}
  \item We release a \textbf{unified civil-law corpus} that spans two disjoint branches of the Polish judiciary in a single coherent package, addressing the strong English-language and common-law tilt of legal-NLP resources.
  \item We provide \textbf{large-scale LLM-based enrichment} of every judgment, with documented prompts, schemas, model version, and per-field statistics, enabling both downstream use and methodological scrutiny.
  \item We treat \textbf{silver-label framing as an explicit design commitment} rather than a post hoc disclaimer, as elaborated in \S\ref{sec:intro:silver}.
  \item We supply a complete \textbf{documentation infrastructure}, including a Datasheet, an Evaluation Card, Croissant and Responsible AI metadata, and reproducibility artifacts that accompany the data release.
  \item We provide a \textbf{descriptive characterization} of the corpora in terms of coverage, structure, citation patterns, and pseudonymization fidelity, enabling cross-branch descriptive comparison between common and administrative courts.
\end{enumerate}

\subsection{Silver labels by design --- not by accident}
\label{sec:intro:silver}

\begin{silverbox}{Design Commitment: Silver Labels by Design}
\textbf{Design commitment.} The analytical fields released with \juddgespl{} are \textbf{silver labels}: produced by a single LLM under uniform prompting, without per-document human review. We make this commitment explicit because the alternative --- calling LLM outputs ``annotations'' without qualification --- has produced a wave of resources whose quality is silently overstated and whose downstream benchmark numbers are not reproducible across model versions.
\end{silverbox}

We adopt this framing as a positive methodological stance rather than a hedge. The analytical fields in \juddgespl{} are described as silver labels throughout the release---never as \emph{annotations} without qualification---and the corpus is offered as a structurally enriched resource rather than as a benchmark with held-out test sets, since meaningful benchmarking would require validated ground truth that we deliberately do not claim to provide. To support informed downstream use we document the error modes that callers should anticipate (\S\ref{sec:pipeline:quality}, \S\ref{sec:limitations:limitations}) and we point to gold-label subsets within the broader JuDDGES family, such as \texttt{pl-swiss-franc-loans}, for tasks that require human-validated labels. Taken together, these choices position \juddgespl{} as a transparent, large-scale silver-label corpus whose intended uses and known limitations are part of the release rather than caveats appended to it.

\subsection{Roadmap}
\label{sec:intro:roadmap}

\S\ref{sec:related-work} reviews related resources; \S\ref{sec:acquisition} describes acquisition and licensing; \S\ref{sec:pipeline} details the enrichment pipeline; \S\ref{sec:characterization} characterizes the corpora; \S\ref{sec:datasheet} presents the Datasheet and Evaluation Card; \S\ref{sec:limitations} enumerates limitations and intended use; \S\ref{sec:release} covers release and infrastructure.


% Related work
% 02-related-work.tex
% Goal: position JuDDGES-pl against (a) existing legal NLP benchmarks
% and corpora, (b) civil-law / non-English legal resources, (c) the
% literature on LLM-as-annotator and silver labels.

\section{Related work}
\label{sec:related-work}

\subsection{Legal NLP benchmarks and resources}
\label{sec:related:benchmarks}

Most established legal NLP benchmarks are dominated by English-language, common-law material. LegalBench~\citep{guhaLegalbenchCollaborativelyBuilt2023a} aggregates 162 collaboratively built tasks with a strong U.S. common-law focus, while LexGLUE~\citep{chalkidis-etal-2022-lexglue} bundles seven tasks drawn primarily from EU institutions and U.S./U.K. courts and offers only limited coverage of Polish material. Task-specific resources such as CaseHOLD~\citep{zheng2021does}, CUAD~\citep{hendrycks2021cuad}, and BillSum~\citep{kornilova2019billsum} have likewise advanced the field within a largely English-language scope. Multi-jurisdictional corpora, including MultiLegalPile~\citep{niklaus2024multilegalpile689gbmultilinguallegal}, Pile of Law~\citep{NEURIPS2022_bc218a0c}, and the LEXTREME benchmark~\citep{niklaus2023lextreme}, broaden coverage across many languages and legal systems but remain Polish-thin in both volume and task design. Against this backdrop, \juddgespl{} is intended not as a competing benchmark but as a Polish-only, structurally enriched complement that adds analytical fields on top of raw judgments.

\subsection{Polish legal NLP}
\label{sec:related:polish}

Resources for Polish legal NLP have grown but remain comparatively sparse. The earlier \juddges{} release at NeurIPS 2024~\citep{augustyniak2024juddges} introduced a bilingual Polish/U.K. judgment dataset and serves as the direct precursor to the present work. Beyond that, the PolEval shared tasks~\citep{poleval2024} have shaped Polish NLP evaluation more generally, although their legal-domain coverage is limited. On the modelling side, Bielik~\citep{bielik2025arxiv} has emerged as an open Polish LLM with associated pre-training and instruction corpora, and the PLLuM consortium~\citep{pllum2025} has organised a national effort to release Polish-language corpora and instruction data. To our knowledge, none of these resources combines two judicial branches under a single uniform LLM-extracted analytical structure at multi-million-judgment scale, which is the gap \juddgespl{} aims to address.

\subsection{LLM-as-annotator, weak supervision, silver labels}
\label{sec:related:llm-annotator}

A growing body of work reports that large language models can match or approach human inter-annotator agreement on many text-classification and annotation tasks~\citep{gilardi2023chatgpt,ziems2024can,pangakis2023automated,chiang2023llmhumaneval}, suggesting practical value for scaling annotation in low-resource domains. At the same time, careful replications and critiques have documented sensitivity to model version, prompt formulation, and distributional bias, urging caution in treating LLM outputs as ground truth~\citep{reiss2023testing,pangakis2023automated}. These methodological concerns sit alongside an older tradition of weak supervision and programmatic labelling, exemplified by Snorkel~\citep{ratner2017snorkel}, which has long argued for explicit treatment of label noise and provenance. Taken together, the positive findings and the documented critiques motivate our decision to frame the LLM-extracted fields in \juddgespl{} as silver labels rather than as gold annotations, and to expose model identity, prompts, and validation samples so that downstream users can reason about label quality.

\subsection{Documentation standards for ML datasets}
\label{sec:related:documentation}

A set of complementary documentation standards has emerged to support responsible release of machine-learning datasets and models. Datasheets for Datasets~\citep{gebru2021datasheets} provide structured questions covering motivation, composition, collection, and recommended uses. Data Cards~\citep{pushkarna2022data} and Model Cards~\citep{mitchell2019model} further standardise human-readable summaries of datasets and trained models, including intended use and known limitations. The Croissant metadata format~\citep{akhtar2024croissant}, together with the NeurIPS Responsible AI extension, adds a machine-readable layer aligned with current submission requirements. \juddgespl{} follows all four of these standards.


% Acquisition & licensing
% 03-acquisition-licensing.tex
% Goal: describe data sources, copyright basis, GDPR / privacy posture,
% and version pinning.

\section{Acquisition and licensing}
\label{sec:acquisition}

\subsection{Source portals}
\label{sec:acquisition:sources}

\juddgespl{} draws on the two official Polish judgment portals that, between them, cover the common and administrative court systems. Common-court rulings are obtained from \url{orzeczenia.ms.gov.pl}, which exposes an official API returning XML payloads; the snapshot used in this release was acquired up to [DATE], with the oldest record dated [DATE]. Administrative-court rulings are obtained from \url{orzeczenia.nsa.gov.pl}, which provides no API and is therefore harvested via HTML scraping under respectful rate limits, together with a re-acquisition policy that revisits records the portal has retroactively edited. The full acquisition procedure, including parser details and incremental update logic, is described in \juddges{} 2024 \S3~\citep{augustyniak2024juddges} and is reused here without modification.

\subsection{Copyright and licensing}
\label{sec:acquisition:copyright}

The authors interpret Art.~4 pkt 2 of the Polish \emph{Ustawa o prawie autorskim i prawach pokrewnych} as placing court judgments (\emph{orzeczenia s\k{a}d\'ow}) within the category of \emph{dokumenty urz\k{e}dowe} and therefore outside copyright protection. On this reading, the raw judgment text redistributed in \juddgespl{} is treated as public-domain material whose use does not require a licence from the issuing courts. The enrichment layer contributed by the authors --- LLM-extracted fields, structured metadata, and the accompanying Croissant descriptor --- is an independent creative and curatorial work and is released under \textbf{CC BY 4.0}. We note that this is the authors' reading of the statute rather than a settled court ruling, and downstream users are encouraged to consult their own counsel for jurisdiction-specific reuse questions.

\subsection{Privacy and GDPR / RODO}
\label{sec:acquisition:privacy}

Pseudonymization is performed at source by the publishing courts before judgments are made available on the official portals: names of natural persons are replaced with initials or codes, and sensitive personal data is redacted in line with Polish judicial publication practice. Names of public officials, judges, and corporate parties are typically retained, as legally permitted in the context of public adjudication. The authors apply no additional de-anonymization step, do not attempt to recover redacted identities, and do not cross-reference judgments with external personal-data sources. Redistributors of \juddgespl{} or derivative artefacts are required to comply with RODO (the Polish implementation of the GDPR) and must refrain from re-identification attempts against the pseudonymized natural persons appearing in the corpus.

\subsection{Versioning}
\label{sec:acquisition:versioning}

The corpus snapshot underlying this release is pinned to [DATE], and subsequent re-acquisitions are published as new snapshots rather than in-place edits. Hugging Face dataset versions correspond directly to commit hashes in the dataset repository, and the Croissant \texttt{version} field tracks these releases so that any cited result can be reproduced against an exact corpus state.


% Enrichment pipeline (methodological core)
% 04-enrichment-pipeline.tex
% Goal: describe the LLM extraction pipeline rigorously enough that it
% is reproducible, including version pinning of the model.
% This is the methodological core of the paper.

\section{Enrichment pipeline}
\label{sec:pipeline}

\subsection{Motivation: from raw text to analysis-ready}
\label{sec:pipeline:motivation}

Raw judgment text is difficult to evaluate over at scale: documents are long, structurally heterogeneous across courts and decades, and written in specialized legal language that resists straightforward parsing. The two prior strategies for converting such text into analysis-ready fields each have well-understood limitations. Human annotation produces high-fidelity labels but does not scale to the millions of judgments that constitute a national case-law corpus, while rule-based extraction is brittle across courts and time as templates, formatting conventions, and citation styles drift. LLM-based extraction with uniform prompting offers a different trade-off, exchanging guaranteed exactness for uniform coverage at corpus scale, with characteristic failure modes that we document rather than dismiss.

\subsection{Model and prompting}
\label{sec:pipeline:model}

We perform enrichment with Google Gemini 2.5 Pro, accessed via the Google AI API, with the version pinned at [exact API version string] so that the released artefacts can be associated with a single model behaviour for reproducibility. Extraction is schema-conditioned, with one prompt per field family rather than a single monolithic call, and each prompt includes few-shot examples drawn from a held-out development set so that no example used to shape the prompts overlaps with released analytical fields. Inference is performed in batched calls with deterministic decoding (\texttt{temperature=0}) where the API supports it, and the full prompt texts are reproduced in Appendix~\ref{sec:appendix:prompts} so that downstream users can inspect, replicate, or critique each extraction step.

\subsection{Extracted fields}
\label{sec:pipeline:fields}

Table~\ref{tab:fields} summarises the fields released in \juddgespl{} and their provenance. The release deliberately separates two tiers: silver-label fields produced by Gemini 2.5 Pro, which encode analytical content such as factual state, legal state, summary, thesis, legal bases, keywords, title, and issue date; and source-portal fields such as court name, judge panel, and signature, which are preserved verbatim from the originating portals. This split lets users decide per-field whether to treat the value as ground truth (for the preserved metadata) or as a model-generated annotation requiring downstream validation.

\begin{table}[h]
\centering
\small
\begin{tabular}{lll}
\toprule
\textbf{Field} & \textbf{Type} & \textbf{Source} \\
\midrule
\texttt{factual\_state} & string & LLM-extracted (Gemini 2.5 Pro) \\
\texttt{legal\_state} & string & LLM-extracted \\
\texttt{extracted\_summary} & string & LLM-extracted \\
\texttt{extracted\_thesis} & string & LLM-extracted \\
\texttt{extracted\_legal\_bases} & list of statute references & LLM-extracted, regex-validated \\
\texttt{extracted\_keywords} & list of strings & LLM-extracted \\
\texttt{extracted\_title} & string & LLM-extracted \\
\texttt{extracted\_date\_issued} & date & LLM-extracted, parsed \\
\texttt{court\_name}, \texttt{judges}, \texttt{signature}, \dots & various & preserved from source portal \\
\bottomrule
\end{tabular}
\caption{Fields released in \juddgespl{}. LLM-extracted fields are silver labels; source-portal fields are preserved verbatim.}
\label{tab:fields}
\end{table}

\subsection{Quality, cost, and runtime}
\label{sec:pipeline:quality}

Aggregate compute, monetary cost, and wall-clock runtime for the full extraction run are reported as [TODO] and [TODO], pending the final accounting tied to the pinned API version. To characterise extraction behaviour qualitatively, a Polish-speaking lawyer inspected a sample of [TODO: $N$] documents and identified recurring error modes, including truncated reasoning sections in long judgments, missed citations to repealed statutes, and over-summarization of short procedural judgments where the model compresses already-terse text. We deliberately do not report a corpus-level aggregate accuracy, and this choice is methodological rather than evasive: there is no validated gold reference at corpus scale against which such a number could be computed, and substituting LLM-vs-LLM agreement as a proxy would conflate fluency-correlation with construct validity while obscuring the silver-label nature of the resource. Reporting a single accuracy figure would, in our view, mislead downstream users more than the qualitative characterisation we provide here and the validated subset described below.

\subsection{Silver labels: what this means in practice}
\label{sec:pipeline:silver}

The methodological choices above carry concrete operational consequences for users of \juddgespl{}. Because the analytical fields are produced uniformly by a single LLM rather than validated against human judgement at corpus scale, they should be treated as a structured layer of model-generated annotations whose utility depends on the downstream task. The callout below states the resulting commitment explicitly, separating uses that are well-supported by silver labels from uses that require gold-standard supervision.

\begin{silverbox}{Silver-Label Commitment: Permitted and Prohibited Uses}
\textbf{Silver-label commitment.} The analytical fields above are produced uniformly by a single LLM and have not been human-validated at corpus scale. Users:
\begin{itemize}
  \item \textbf{may} use them for exploratory analysis, descriptive statistics, model pretraining, weak supervision, retrieval indexing;
  \item \textbf{must not} use them as gold labels for benchmarking, regulatory claims, or downstream legal decisions;
  \item \textbf{must} cite the model version when reporting analyses, since LLM behavior shifts across model releases.
\end{itemize}
A separate validated subset (\texttt{JuDDGES/pl-swiss-franc-loans}) provides gold annotations for tasks requiring high-fidelity labels.
\end{silverbox}


% Corpus characterization
% 05-corpus-characterization.tex
% Goal: descriptive statistics that justify the claim
% "evaluation-relevant" without running models.

\section{Corpus characterization}
\label{sec:characterization}

\subsection{Coverage}
\label{sec:characterization:coverage}

\juddgespl{} comprises \textbf{2{,}691{,}842 judgments} merged from the two source subsets (Table~\ref{tab:corpus-overview}): the common-courts subset (\texttt{pl-court}, 437{,}450 judgments published on \url{orzeczenia.ms.gov.pl}) and the administrative-courts subset (\texttt{pl-nsa}, 2{,}254{,}392 judgments from \url{orzeczenia.nsa.gov.pl}). All rows are released as a single \texttt{train} split with a unified 53-column schema; the \texttt{source\_dataset} column distinguishes the two subsets.

\begin{table}[h]
\centering
\small
\begin{tabular}{lrrr}
\toprule
\textbf{Subset} & \textbf{Judgments} & \textbf{Schema columns} & \textbf{LLM-extracted columns} \\
\midrule
\texttt{pl-court} (common courts)         & 437{,}450    & 53 & 12 \\
\texttt{pl-nsa} (administrative courts)   & 2{,}254{,}392 & 53 & 12 \\
\midrule
\textbf{\juddgespl{} (combined)}          & \textbf{2{,}691{,}842} & \textbf{53} & \textbf{12} \\
\bottomrule
\end{tabular}
\caption{\juddgespl{} corpus overview. Both subsets are normalized to a single 53-column schema; 12 columns carry LLM-extracted silver-label content and the remainder are preserved from the source portals or are provenance fields. The \texttt{extracted\_tax\_specific} field present in the source datasets was dropped during the merge ($<\!0.2\%$ coverage).}
\label{tab:corpus-overview}
\end{table}

To complement the aggregate counts in Table~\ref{tab:corpus-overview}, we report three breakdowns of the \juddgespl{} corpus in Appendix~\ref{sec:appendix:stats}: the distribution of judgments across court levels (Supreme Court, appellate, regional, and district within the common-courts subset; NSA and WSA within the administrative subset), a yearly timeline of judgment counts spanning the full coverage window of both source portals, and a breakdown by case type and procedural category. These distributions characterize how the corpus is concentrated across the Polish judiciary and provide the contextual axes against which all subsequent analyses in this section should be read.

\subsection{Document structure}
\label{sec:characterization:structure}

We characterize the structural composition of \juddgespl{} along three axes that are reported in full in Appendix~\ref{sec:appendix:stats}. First, we provide length distributions in tokens and in number of structural sections, computed separately for the common and administrative subsets, since the two portals expose different rendering conventions for judgment text. Second, we report section-presence statistics for the four canonical components of a Polish judgment, namely the operative part (\emph{sentencja}), the reasoning (\emph{uzasadnienie}), dissenting opinions, and \emph{glosa} commentary, with each component tracked as a binary indicator over judgments. Third, we report the ratio of \emph{sentencja} length to \emph{uzasadnienie} length, which captures the balance between the operative and the explanatory portions of a judgment and varies systematically by court level.

\subsection{Citation patterns}
\label{sec:characterization:citations}

We provide three citation-pattern analyses that characterize how judgments in \juddgespl{} reference legal authority. First, we release a statute-citation graph derived from the \texttt{extracted\_legal\_bases} field, together with a ranking of the most frequently cited statutes; this view reflects the Polish civil-law tradition in which statutory provisions are the dominant citation target. Second, we quantify case-citation patterns; case citations are intrinsically rare in a civil-law jurisdiction and we report their incidence as such, rather than treating them as a primary signal. Third, we report cross-branch citation behavior between the common and administrative subsets as a sanity check: administrative-court judgments rarely cite common-court judgments, and we quantify this rate to document the structural separation between the two branches rather than to motivate a downstream task. Numerical values for all three analyses are deferred to Appendix~\ref{sec:appendix:stats}.

\subsection{Extraction-field statistics}
\label{sec:characterization:extraction-stats}

\paragraph{Coverage per LLM-extracted field.}
Table~\ref{tab:extraction-coverage} reports the fraction of judgments for which each LLM-extracted field carries a non-empty value, computed across all 2{,}691{,}842 judgments in the released corpus. Coverage is bimodal: a single high-coverage field, \texttt{extracted\_legal\_bases}, is populated for 81.3\% of judgments because every judgment that cites at least one statute yields a non-empty list; in contrast, the free-text analytical fields (\texttt{extracted\_summary}, \texttt{extracted\_thesis}, \texttt{extracted\_factual\_state}, \dots) populate at 12--16\% combined coverage, reflecting the proportion of judgments for which the source-portal text contains a passage the extractor could ground its output on. Per-subset coverage is broadly similar, with the common-courts subset (\texttt{pl-court}) consistently 2--3 percentage points above the administrative subset for analytical fields, and 10 percentage points higher for \texttt{extracted\_legal\_bases}; we attribute this to richer reasoning-section content (\emph{uzasadnienie}) in common-court judgments. Two long-tail fields, \texttt{extracted\_legal\_analysis} and \texttt{extracted\_judgment\_specific}, drop to single-digit percentages on the administrative subset and are intended primarily for downstream filtering rather than as dense supervision signal.

\begin{table}[h]
\centering
\small
\begin{tabular}{lrrr}
\toprule
\textbf{Field} & \textbf{\texttt{pl-nsa}} & \textbf{\texttt{pl-court}} & \textbf{Combined} \\
\midrule
\texttt{extracted\_legal\_bases}        & 79.7\% & 89.6\% & \textbf{81.3\%} \\
\texttt{extracted\_summary}             & 15.2\% & 17.6\% & 15.6\% \\
\texttt{extracted\_factual\_state}      & 14.8\% & 17.0\% & 15.1\% \\
\texttt{extracted\_thesis}              & 13.7\% & 16.6\% & 14.1\% \\
\texttt{extracted\_legal\_state}        & 12.8\% & 15.9\% & 13.3\% \\
\texttt{extracted\_outcome}             & 13.2\% & 13.0\% & 13.1\% \\
\texttt{extracted\_parties}             & 12.7\% & 12.0\% & 12.6\% \\
\texttt{extracted\_keywords}            & 12.4\% & 12.1\% & 12.4\% \\
\texttt{extracted\_legal\_references}   & 12.2\% & 11.5\% & 12.1\% \\
\texttt{extracted\_legal\_concepts}     & 12.1\% & 11.3\% & 12.0\% \\
\texttt{extracted\_legal\_analysis}     &  8.5\% &  3.6\% &  7.7\% \\
\texttt{extracted\_judgment\_specific}  &  7.3\% &  3.8\% &  6.8\% \\
\bottomrule
\end{tabular}
\caption{Coverage of LLM-extracted fields in \juddgespl{}. Coverage is the fraction of judgments where the field is non-null and non-empty (for list-valued fields, contains at least one element). Computed exhaustively over all 2{,}691{,}842 judgments. The \texttt{extracted\_tax\_specific} field present in the source datasets was dropped during the merge ($<\!0.2\%$ coverage).}
\label{tab:extraction-coverage}
\end{table}

Beyond coverage, we release per-field length distributions for each of the twelve LLM-extracted fields and deduplication-aware diversity metrics for \texttt{extracted\_keywords} that account for near-duplicate keyword surface forms across judgments. These analyses are released alongside the corpus and reported in Appendix~\ref{sec:appendix:stats}.

\subsection{Pseudonymization fidelity}
\label{sec:characterization:pseudonymization}

We provide an empirical audit of pseudonymization fidelity in \juddgespl{} on a stratified sample of [TODO: $N$] judgments drawn proportionally across court level, year, and source subset. The audit measures the residual presence of personal-identifier patterns that should have been removed by the source portals before publication, including PESEL national-identification numbers detected by regular expression, postal-address patterns, and full personal names. We further report how the residual-identifier rate varies across courts and across years, since pseudonymization conventions on the source portals have evolved over time and differ between common and administrative courts. We emphasize that this audit is an empirical characterization of the released text and is explicitly not a privacy guarantee; users handling \juddgespl{} for downstream applications should perform their own risk assessment in accordance with the data-protection requirements of their jurisdiction.

\subsection{Cross-branch comparison}
\label{sec:characterization:cross-branch}

To document how the common-courts and administrative-courts subsets relate to each other, we report descriptive cross-branch comparisons along three axes. We compute token-level vocabulary overlap between the two subsets, which characterizes lexical divergence between common and administrative judicial language. We compute statute-citation overlap from the \texttt{extracted\_legal\_bases} field, identifying the set of statutes cited in both subsets versus those cited exclusively in one. We further report topic coverage by case type, comparing the procedural categories that are well represented in each subset. These descriptive comparisons are intended to characterize the joint corpus rather than to assert any particular cross-branch transfer or generalization claim, and the corresponding tables are reported in Appendix~\ref{sec:appendix:stats}.


% Datasheet and Evaluation Card
% 06-datasheet-evaluation-card.tex
% Goal: condensed Datasheet [Gebru et al. 2021] + an Evaluation Card
% declaring what evaluative claims this resource is designed to support.

\section{Datasheet and Evaluation Card}
\label{sec:datasheet}

\subsection{Datasheet}
\label{sec:datasheet:datasheet}

We accompany \juddgespl{} with a condensed Datasheet following the template of \citet{gebru2021datasheets}, covering motivation, composition, collection, preprocessing, intended uses, distribution, and maintenance. The body of the paper has already addressed each of these dimensions in passing: motivation and composition are introduced in the corpus overview, collection and preprocessing are described alongside the silver-label pipeline, and distribution and maintenance are detailed in the release section. To keep this section short while remaining standards-compliant, the full Datasheet template, with one paragraph per question, is provided in Appendix~\ref{sec:appendix:datasheet}. Readers seeking question-by-question answers, including provenance, sensitive content, and known limitations, should consult that appendix.

\subsection{Evaluation Card}
\label{sec:datasheet:eval-card}

Beyond the Datasheet, we provide an Evaluation Card (Table~\ref{tab:eval-card}) that explicitly declares which classes of evaluative claim \juddgespl{} is designed to support and which it is not. The purpose of such a card is to set expectations up front: a corpus built primarily with silver labels can legitimately ground some research uses while being inappropriate for others, and conflating the two is a recurring failure mode in legal NLP releases. We therefore enumerate seven claim types and mark each as either supported or unsupported, with a short justification, so that downstream users can match their research questions to the resource's evidentiary basis before running experiments.

\begin{table}[h]
\centering
\small
\begin{tabular}{lcl}
\toprule
\textbf{Claim type} & \textbf{Supported?} & \textbf{Notes} \\
\midrule
Pretraining / continued pretraining of legal LMs in Polish & Yes & Silver labels acceptable; raw text high-fidelity. \\
Weak-supervision training of extraction models & Yes & Silver labels are the intended use. \\
Descriptive / longitudinal legal analytics & Yes & Acknowledge model-version drift. \\
Cross-branch generalization studies & Yes & Two disjoint branches enable this. \\
\textbf{Benchmarking} with held-out test sets on LLM-extracted fields & No & Use validated subsets instead. \\
Privacy guarantees beyond court-applied pseudonymization & No & Document anonymization is heterogeneous. \\
Individual legal decision-making & No & Out of scope; not legal advice. \\
\bottomrule
\end{tabular}
\caption{Evaluation Card: claims this resource is designed to support.}
\label{tab:eval-card}
\end{table}

The four supported uses follow directly from the corpus design. Pretraining and continued pretraining of Polish legal language models are core uses because the underlying judgment text is high-fidelity and silver labels are not on the loss path, so extraction noise does not propagate into the model. Weak-supervision training of extraction models is, by construction, the intended use: silver labels are produced at scale precisely so that downstream extractors can be trained or distilled from them, with the understanding that any reported held-out numbers must come from a validated subset rather than from the silver labels themselves. Descriptive and longitudinal legal analytics are likewise supported, provided that users report which model version produced the labels and acknowledge the model-version drift this introduces over time. Finally, the resource supports descriptive cross-branch comparison between common and administrative courts; cross-branch generalization claims, in the strict sense, would require gold-label validation we do not provide, and we encourage users to frame branch-level findings descriptively.

The three unsupported uses are excluded for equally concrete reasons. Benchmarking with held-out test sets defined directly on LLM-extracted fields is not supported because there is no validated gold reference for those fields at corpus scale, so reported scores would conflate extractor agreement with task performance; users should benchmark on the validated subsets instead. We do not claim privacy guarantees beyond the pseudonymization already applied by the courts at publication time, since anonymization practice is heterogeneous across chambers and years and we do not re-anonymize. Individual legal decision-making is explicitly out of scope: \juddgespl{} is a research resource, not legal advice, and must not be used to determine outcomes for specific persons or cases.

\subsection{Responsible AI metadata}
\label{sec:datasheet:rai}

All Croissant Responsible AI fields are populated for the release, including data-collection context, sensitive-content flags, intended and out-of-scope uses, and maintenance contact, in line with the Evaluation Card above. The full machine-readable metadata is shipped as \texttt{juddges-pl.json} alongside the corpus; see \S\ref{sec:release} for the distribution channels and the resolved file URL.


% Limitations and intended use
% 07-limitations-intended-use.tex
% This section is the silver-label commitment in operational form.
% Reviewers and downstream users should read it before depending on
% the resource.

\section{Limitations and intended use}
\label{sec:limitations}

\begin{silverbox}{Reviewer and User Notice: Silver-Label Commitment in Operational Form}
\textbf{Reviewer / user notice.} This section is the silver-label commitment in operational form. Reviewers and downstream users should read it before depending on the resource.
\end{silverbox}

\subsection{Limitations}
\label{sec:limitations:limitations}

\juddgespl{} is scoped to a single jurisdiction, namely Poland, and we make no cross-jurisdictional claims without the addition of further resources. The analytical fields produced by our extraction pipeline are silver labels rather than human-validated annotations: they are not verified at corpus scale, they are subject to model-version drift, and we therefore intentionally refrain from reporting an aggregate accuracy figure (see \S\ref{sec:pipeline:quality}). Court-applied pseudonymization is heterogeneous across courts and across time, and consequently does not amount to a uniform privacy guarantee for downstream users. Because only published judgments enter the corpus, and because publisher selection is non-random, the resulting sample carries publication bias that should be acknowledged in any descriptive analysis. Temporal coverage is denser in recent years and conspicuously sparser before 2000, which constrains long-horizon longitudinal studies. Finally, the corpus is monolingual in Polish, and we deliberately do not include machine translations into other languages.

\subsection{Intended uses}
\label{sec:limitations:uses}

We envisage \juddgespl{} as a substrate for pretraining and weak-supervision pipelines, for retrieval indexing over Polish case law, and for descriptive analytics over the published judiciary record. The resource also supports fairness audits, broader civil-law NLP research, and civic-tech tooling, in each case under appropriate disclaimers regarding the silver-label nature of the analytical fields. These use cases align with the operational commitments stated in \S\ref{sec:limitations:limitations} and with the model-version disclosure requirement in \S\ref{sec:limitations:model-version}.

\subsection{Out-of-scope uses}
\label{sec:limitations:out-of-scope}

Users should not benchmark systems against the LLM-derived analytical fields without first conducting their own independent validation against gold annotations on the slice of interest. Privacy-sensitive applications should not rely on the corpus without performing additional anonymization beyond what individual courts already apply. The resource is also not intended for individual legal advice, for regulatory decisions, or for any other high-stakes individual outcome, and downstream tooling should make this scope explicit to its end users.

\subsection{Model-version disclosure}
\label{sec:limitations:model-version}

Because LLM outputs are model-specific, users analyzing the extracted fields should cite both the Gemini model version recorded in \S\ref{sec:pipeline:model} and the dataset version they consumed. This pairing is what makes results derived from the silver labels reproducible and comparable across studies.


% Release and infrastructure
% 08-release-infrastructure.tex

\section{Release and infrastructure}
\label{sec:release}

\subsection{Hugging Face release}
\label{sec:release:hf}

The canonical artifact accompanying this paper is \texttt{JuDDGES/juddges-pl}, a single unified release on the Hugging Face Hub that exposes two configurations, \texttt{pl-court} and \texttt{pl-nsa}, corresponding to the common-court and administrative-court subcorpora of \juddgespl{}. The release is publicly accessible without gating, authentication, or click-through agreements, so that any researcher can stream or download the data directly. Wroclaw University of Science and Technology (WUST) commits to maintaining the dataset on Hugging Face under best-effort institutional support; the maintenance plan is documented in the Croissant \texttt{dataReleaseMaintenancePlan} field. The earlier component datasets \texttt{JuDDGES/pl-court-raw-enriched} and \texttt{JuDDGES/pl-nsa-enriched} served as upstream sources during pipeline development and remain available for provenance, but the canonical release for this paper is \texttt{JuDDGES/juddges-pl}.

\subsection{Croissant metadata}
\label{sec:release:croissant}

We ship a validated \texttt{juddges-pl.json} Croissant descriptor for the unified release alongside per-component fallbacks for the two configurations, so that downstream consumers can resolve metadata at either granularity. The descriptor populates the full set of Responsible AI fields recommended by the Croissant RAI vocabulary, namely \texttt{dataCollection}, \texttt{dataCollectionType}, \texttt{dataCollectionRawData}, \texttt{dataPreprocessingProtocol}, \texttt{dataAnnotationProtocol}, \texttt{dataAnnotationPlatform}, \texttt{dataAnnotationAnalysis}, \texttt{annotationsPerItem}, \texttt{annotatorDemographics}, \texttt{machineAnnotationTools}, \texttt{dataUseCases}, \texttt{dataBiases}, \texttt{personalSensitiveInformation}, \texttt{dataLimitations}, \texttt{dataReleaseMaintenancePlan}, and \texttt{dataSocialImpact}. Together these fields document acquisition provenance, preprocessing and annotation protocols, the role of LLM-based silver labelling, intended uses and known biases, the treatment of personal data, and the maintenance plan referenced in Section~\ref{sec:release:hf}.

\subsection{Reproducibility artifacts}
\label{sec:release:repro}

Beyond the dataset itself, we release the artifacts needed to reproduce the enrichment layer end-to-end. The extraction prompts and per-field schemas used for silver labelling are reproduced in Appendix~\ref{sec:appendix:prompts}, and the Croissant descriptors are regenerated by the script \texttt{croissant/generate\_croissant.py} bundled with the codebase. Exploratory data-analysis notebooks that produced the corpus statistics in this paper are described in Appendix~\ref{sec:appendix:samples} and shipped in the supplementary ZIP.

\subsection{License}
\label{sec:release:license}

The enrichment layer of \juddgespl{}, comprising the silver labels, derived fields, and accompanying metadata, is released under the Creative Commons Attribution 4.0 (CC BY 4.0) license. The raw judgment text is in the public domain under Polish copyright law, as discussed in \S\ref{sec:acquisition:copyright}.


% Conclusion
% 09-conclusion.tex

\section{Conclusion}
\label{sec:conclusion}

We presented \juddgespl{}, a large-scale enriched Polish judgment corpus addressing the civil-law gap in legal NLP resources. By combining two disjoint judicial branches, applying a uniform LLM-based enrichment pipeline, and committing to \textbf{silver-label framing} as a first-class design choice, we offer a resource whose evaluative role is documented as carefully as its content. We invite the community to build gold-label subsets, validation studies, and downstream benchmarks atop this foundation.


% Bibliography
\bibliography{bib}

%%%%%%%%%%%%%%%%%%%%%%%%%%%%%%%%%%%%%%%%%%%%%%%%%%%%%%%%%%%%
\appendix
% appendix.tex
% Appendices A-E. \appendix is invoked once in main.tex before this file.

\section{Extraction prompts and field schemas}
\label{sec:appendix:prompts}

For each of \texttt{factual\_state}, \texttt{legal\_state}, \texttt{extracted\_summary}, \texttt{extracted\_thesis}, \texttt{extracted\_legal\_bases}, \texttt{extracted\_keywords}, \texttt{extracted\_title}, \texttt{extracted\_date\_issued}: full prompt text, few-shot examples (drawn from held-out development set), and parser/post-processor.

% TODO: import content from appendix_extraction.tex once polished.

\section{Datasheet for Datasets (full template)}
\label{sec:appendix:datasheet}

Following Gebru et al.~\citep{gebru2021datasheets}: motivation, composition, collection process, preprocessing/cleaning/labeling, uses, distribution, maintenance.

% TODO: port the full datasheet questionnaire.

\section{Sample records}
\label{sec:appendix:samples}

% TODO: include 2-3 fully-anonymized example records per config, showing
% all fields side-by-side. Use the examplebox tcolorbox defined in main.tex.

\section{Detailed statistics tables}
\label{sec:appendix:stats}

% TODO: full tables for the figures referenced in
% Section~\ref{sec:characterization}.

\section{NeurIPS 2026 Paper Checklist}
\label{sec:appendix:checklist}

The mandatory NeurIPS 2026 paper checklist is included separately at the end of the PDF (after this appendix). See \texttt{sections/checklist.tex}; the checklist file is invoked from \texttt{main.tex}.


%%%%%%%%%%%%%%%%%%%%%%%%%%%%%%%%%%%%%%%%%%%%%%%%%%%%%%%%%%%%
\newpage
% Mandatory NeurIPS 2026 paper checklist --- desk reject if missing.
\input{sections/checklist}

\end{document}
